\section{从感知机到循环神经网络}

\subsection{回顾感知机}

在上一讲中，我们学习了感知机的基本结构：给定输入 $x^{(1)}, x^{(2)}, \ldots, x^{(m)}$，每个输入乘以对应权重 $w_1, w_2, \ldots, w_m$，求和后通过非线性激活函数产生输出。

\begin{figure}[H]
\centering
\includegraphics[width=0.7\textwidth]{slides-images/slide-10.jpg}
\caption{感知机结构回顾\protect\footnotemark}
\end{figure}
\footnotetext{Slides 第10页。}

用数学公式表示：

$$z = \sum_{i=1}^{m} w_i x^{(i)} + b$$

$$\hat{y} = g(z)$$

其中：
\begin{itemize}
    \item $x^{(i)}$：第 $i$ 个输入特征
    \item $w_i$：对应的权重
    \item $b$：偏置项
    \item $g(\cdot)$：非线性激活函数
    \item $\hat{y}$：预测输出
\end{itemize}

\subsection{感知机处理序列数据的局限}

将感知机（或前馈网络）直接应用于序列数据时，我们遇到了一个根本性问题：\textbf{每个时间步的预测完全独立}，没有利用其他时间步的信息。

\begin{figure}[H]
\centering
\includegraphics[width=0.75\textwidth]{slides-images/slide-13.jpg}
\caption{前馈网络独立处理每个时间步：无法利用历史信息\protect\footnotemark}
\end{figure}
\footnotetext{Slides 第13页。}

具体来说，如果我们用前馈网络独立处理每个时间步：

$$\hat{y}_t = f(x_t)$$

那么时间步 $t=2$ 的预测 $\hat{y}_2$ 仅仅取决于输入 $x_2$，完全忽略了 $x_0$ 和 $x_1$ 中可能包含的关键信息。

这就好比只看球的当前位置而不考虑历史轨迹——预测必然是盲目的。

\subsection{引入循环：Recurrence 的核心思想}

为了解决这一问题，我们需要一种机制将过去时间步的信息传递给当前时间步。核心思想是：

\begin{importantbox}{循环神经网络的核心：隐藏状态}
定义一个\textbf{内部状态（Hidden State）}$h_t$，它在时间步之间传递，携带了网络对过去输入的``记忆''。每个时间步的预测不仅依赖当前输入 $x_t$，还依赖上一步的隐藏状态 $h_{t-1}$：

$$\hat{y}_t = f(x_t, h_{t-1})$$
\end{importantbox}

\begin{figure}[H]
\centering
\includegraphics[width=0.85\textwidth]{slides-images/slide-15.jpg}
\caption{循环神经元：输出依赖当前输入与过去记忆\protect\footnotemark}
\end{figure}
\footnotetext{Slides 第15页。}

上图左侧展示了 RNN 的折叠视图（Folded View），右侧展示了展开视图（Unrolled View）。隐藏状态 $h_t$ 像一条链条一样在时间步之间传递，使得网络能够``记住''之前看到的信息。

\subsection{本章小结}

从感知机到 RNN 的演进思路清晰：前馈网络缺乏处理序列的能力，因此我们引入隐藏状态作为时间步之间的信息桥梁，形成了循环（Recurrence）的概念。这是理解所有序列模型的基础。

\newpage
